%! TEX root = produk_kota.tex
% the main TeX file which is intended to compile, :VimtexReload after adjustment
   % this file should be included in the main file
% :h vimtex-tex-root
\documentclass[../rumukota.tex]{subfiles}
%ensure the class of docs

\begin{document}

\section{Produk dan Unsur-Unsur Rancang Kota}

\subsection{Pendahuluan}
Modul ini disusun berdasarkan subpokok bahasan 1.2 yang memfokuskan pada pemahaman fundamental mengenai \textit{urban design} atau rancang kota melalui dua pilar utama: produk-produk yang dihasilkan dan unsur-unsur pembentuknya. Rancang kota bukan sekadar disiplin estetika, melainkan sebuah proses sistematis yang menghasilkan instrumen-instrumen kebijakan dan rencana fisik yang terukur.

Tujuan utama dari bab ini adalah memberikan pemahaman mendalam bagi mahasiswa mengenai hierarki produk rancang kota, mulai dari kebijakan makro hingga desain proyek mikro, serta mengenali komponen-komponen fisik yang membentuk lingkungan perkotaan yang harmonis. Seluruh narasi dalam bab ini merujuk pada materi presentasi yang mencakup produk, skala, jenis proyek, produk khusus, hingga representasi visual dalam rancang kota.

\subsection{Produk-Produk Rancang Kota}
Dalam praktik profesional, produk rancang kota tidak bersifat tunggal. Produk tersebut terbagi ke dalam empat kategori utama yang saling berkaitan dan membentuk satu kesatuan alur kerja dari tahap gagasan hingga implementasi di lapangan.

\subsubsection{Design Policy (Kebijakan Rancang)}
\textit{Design policy} merupakan instrumen yang berada pada tingkatan paling awal dan bersifat strategis. Produk ini didefinisikan sebagai metoda perancangan tak langsung (\textit{indirect design method}). Alih-alih langsung memberikan bentuk fisik, kebijakan ini bekerja melalui:
\begin{itemize}
	\item \textbf{Instrumen Peraturan:} Menyediakan landasan hukum atau regulasi yang menjadi syarat mutlak pelaksanaan pembangunan di suatu wilayah.
	\item \textbf{Program Investasi:} Mengintegrasikan aspek ekonomi dengan rancangan, sehingga rencana yang dibuat memiliki kepastian sumber daya dan instrumen pendukung lainnya yang menyebabkan rancangan tersebut dapat dilaksanakan secara operasional.
	\item \textbf{Kerangka Kerja:} Bertindak sebagai penuntun (\textit{framework}) bagi keseluruhan proses perancangan dari awal hingga akhir.
	\item \textbf{Prinsip Keluwesan:} Kebijakan rancang harus dirumuskan dengan tingkat fleksibilitas atau keluwesan yang cukup. Hal ini krusial agar ketika rencana tersebut diturunkan ke skala yang lebih spesifik di tingkat bawahnya, rancangan tersebut tetap dapat diimplementasikan tanpa hambatan birokrasi yang kaku.
\end{itemize}

\subsubsection{Design Plan (Rencana Rancang)}
Setelah kebijakan ditetapkan, langkah selanjutnya adalah visualisasi melalui \textit{design plan}. Produk ini merupakan bentuk penggambaran tiga dimensi (3D) dari seluruh kebijaksanaan rancangan yang telah ditetapkan sebelumnya. Ciri khas dari \textit{design plan} adalah:
\begin{itemize}
	\item \textbf{Wawasan Produk:} Menitikberatkan pada hasil akhir fisik yang ingin dicapai (bentuk bangunan, ruang terbuka, dan infrastruktur).
	\item \textbf{Wawasan Proses:} Mempertimbangkan tahapan-tahapan bagaimana produk tersebut akan terwujud seiring berjalannya waktu.
\end{itemize}

\subsubsection{Design Guidelines (Panduan Rancang)}
\textit{Design guidelines} berfungsi sebagai dokumen teknis yang menjabarkan kebijaksanaan dan rencana ke dalam panduan yang sangat rinci. Fokus utama produk ini adalah untuk menjamin kualitas rancangan pada skala mikro. Beberapa karakteristik penting dari panduan ini antara lain:
\begin{itemize}
	\item \textbf{Sifat Non-Restriktif:} Panduan ini tidak dimaksudkan sebagai kendali atau pembatas yang mengekang kreativitas secara kaku. Sebaliknya, ia disusun untuk mengembangkan kerangka rancangan pada berbagai tingkat, mulai dari kawasan, koridor jalan, tingkat proyek, hingga ke unit terkecil yaitu persil lahan.
	\item \textbf{Penyajian Alternatif:} Panduan rancang kota yang baik harus mampu menyajikan berbagai alternatif bentuk atau pendekatan untuk unsur rancang kota yang spesifik, seperti pengaturan tata letak \textit{plaza}, kompleks perumahan, hingga detail lanskap.
\end{itemize}

\subsubsection{Design Program (Program Rancang)}
\textit{Design program} merupakan wujud nyata dari proses pelaksanaan. Produk ini mengatur seluruh proses perancangan dengan tujuan ganda:
\begin{itemize}
	\item \textbf{Konservasi:} Memelihara dan melestarikan karakter lingkungan fisik yang sudah ada (\textit{existing}).
	\item \textbf{Kreasi:} Menciptakan lingkungan fisik yang sama sekali baru sesuai dengan visi perancangan yang telah ditetapkan.
\end{itemize}

\subsection{Skala Rencana Rancang Kota}
Perancangan kota beroperasi pada berbagai skala geografis dan administratif yang berbeda. Pemahaman terhadap skala ini menentukan tingkat kedetailan dan otoritas yang terlibat.

\subsubsection{City Design (Rancangan Skala Kota)}
Ini adalah tingkatan tertinggi yang mencakup wilayah administratif kota secara utuh. Fokus utamanya adalah:
\begin{itemize}
	\item Menyediakan panduan yang didasarkan pada kepentingan masyarakat luas.
	\item Mempertimbangkan dan mengarahkan substansi perencanaan kota secara makro.
	\item Memberikan konteks rancangan fisik yang menjadi dasar bagi pengambilan keputusan perencanaan serta koordinasi pembangunan, baik yang dilakukan oleh pemerintah maupun sektor swasta.
\end{itemize}

\subsubsection{District Design (Rancangan Skala Kawasan)}
Pada skala ini, perancangan memusat pada area tertentu yang ditentukan berdasarkan fungsi dominan (seperti kawasan perdagangan atau retail) atau berdasarkan pengelompokan lingkungan geografis (seperti kawasan perumahan). Peran utamanya adalah:
\begin{itemize}
	\item Menyediakan prosedur perancangan dan teknik spasial untuk melindungi serta memperkaya aset lingkungan, baik yang alami maupun buatan.
	\item Bertindak sebagai pemicu bagi proyek-proyek khusus yang akan dilaksanakan oleh pemerintah atau swasta.
\end{itemize}

\subsubsection{Project Design (Rancangan Skala Proyek)}
Merupakan skala paling teknis dan spesifik yang berkaitan langsung dengan penanganan tapak (\textit{site}). Karakteristiknya meliputi:
\begin{itemize}
	\item Pengembangan alat kendali untuk memastikan sasaran skala kota dan kawasan tercapai.
	\item Penyeimbangan antara aspek visual dan aspek fungsional secara detail.
	\item Output berupa studi kelayakan (\textit{feasibility study}), paket insentif pembangunan, skema tahapan pembangunan kawasan (misal: koridor jalan arteri), hingga paket renovasi skala kecil.
\end{itemize}

\subsection{Jenis Proyek dalam Rancang Kota}
Lingkup kerja rancang kota mencakup berbagai tipologi proyek yang luas, di antaranya:
\begin{enumerate}
	\item \textit{Infrastructure design}: Perancangan infrastruktur kota.
	\item \textit{New town design}: Perancangan kota baru.
	\item \textit{Large project design}: Perancangan proyek-proyek berskala besar.
	\item \textit{Precinct within neighborhood}: Perancangan zona/wilayah di dalam lingkungan permukiman.
	\item \textit{Urban renewal}: Program peremajaan kawasan kota yang telah mengalami degradasi.
	\item \textit{Sub urban development}: Pengembangan area pinggiran kota.
	\item \textit{Campus}: Mencakup kompleks pendidikan, perkantoran, pusat medis (\textit{medical centre}), dan lain-lain.
\end{enumerate}

\subsection{Produk Khusus Rancang Kota}
Secara legal dan administratif di Indonesia, produk rancang kota bermanifestasi ke dalam beberapa instrumen khusus, yaitu:
\begin{itemize}
	\item \textbf{RTBL (Rencana Tata Bangunan dan Lingkungan):} Dokumen induk yang mengatur penataan bangunan dan lingkungan.
	\item \textbf{Program Bangunan dan Lingkungan:} Rencana aksi pelaksanaan pembangunan fisik.
	\item \textbf{Program Investasi dan Insentif Pengembangan:} Skema ekonomi untuk menarik keterlibatan sektor swasta.
	\item \textbf{Rencana Umum Pengendalian Wujud Bangunan:} Setara dengan \textit{design plan} untuk mengontrol estetika dan massa bangunan.
	\item \textbf{Rencana Detail:} Merupakan bentuk \textit{design guidelines} untuk panduan teknis di lapangan.
	\item \textbf{Administrasi Pengendalian:} Pedoman tata usaha dan prosedur (\textit{administration guidelines}).
	\item \textbf{Arahan Pengendalian Pelaksanaan:} Panduan konkret saat proses konstruksi dan pengembangan berlangsung (\textit{development guidelines}).
\end{itemize}

\subsection{Unsur-Unsur Rancang Kota}
Unsur-unsur ini merupakan elemen fisik dan non-fisik yang menjadi objek utama dalam setiap produk perancangan kota:
\begin{itemize}
	\item \textbf{Land Use:} Tata guna lahan yang menentukan fungsi suatu ruang.
	\item \textbf{Massing and Building:} Gubahan massa dan bentuk bangunan yang membentuk karakter visual kawasan.
	\item \textbf{Circulation and Parking:} Sistem pergerakan kendaraan dan penyediaan fasilitas parkir.
	\item \textbf{Open Space:} Ruang terbuka yang berfungsi sebagai paru-paru dan ruang sosial kota.
	\item \textbf{Pedestrian Way:} Jalur khusus pejalan kaki yang aman dan nyaman.
	\item \textbf{Activity Supports:} Pendukung aktivitas yang memicu keramaian dan keberlangsungan fungsi ruang.
	\item \textbf{Signage:} Sistem penandaan dan informasi visual di ruang publik.
	\item \textbf{Preservation:} Upaya pelestarian bangunan atau kawasan bersejarah.
	\item \textbf{Public Arts:} Integrasi karya seni dalam ruang kota untuk meningkatkan kualitas estetika.
	\item \textbf{Street Furniture:} Perabot jalan (bangku, lampu, tempat sampah) yang menunjang kenyamanan pengguna.
	\item \textbf{View:} Perlindungan terhadap sudut pandang atau pemandangan strategis kota.
\end{itemize}

\subsection{Representasi Visual dalam Rencana Rancang Kota}
Untuk mengomunikasikan rencana kepada \textit{stakeholder}, diperlukan berbagai format representasi visual yang mencakup:
\begin{itemize}
	\item \textbf{Zoning 3D:} Visualisasi pembagian zona (hunian, retail, parkir) dalam bentuk tiga dimensi.
	\item \textbf{Section:} Gambar potongan untuk memperlihatkan hubungan ketinggian bangunan dan ruang jalan.
	\item \textbf{Circulation Plan:} Peta khusus yang mengatur alur pergerakan di dalam kawasan.
	\item \textbf{Parking Plan:} Detail penempatan dan kapasitas fasilitas parkir.
	\item \textbf{Open Space Plan:} Rencana distribusi dan jenis ruang terbuka.
	\item \textbf{Massing \& Building Plan:} Representasi gubahan massa bangunan secara kolektif.
	\item \textbf{Streetscape:} Gambaran rupa wajah jalan dan interaksinya dengan bangunan di tepinya.
	\item \textbf{Signage Plan:} Peta lokasi dan desain elemen penanda kota.
	\item \textbf{Illustrative Design Plan:} Rencana desain yang disajikan secara ilustratif dan komprehensif untuk memudahkan pemahaman visi desain.
	\item \textbf{Perspective (Bird's Eye View):} Sudut pandang mata burung untuk melihat keseluruhan korelasi antarelemen dalam satu kawasan luas.
\end{itemize}

\bibliographystyle{apalike}
\bibliography{../filename.bib} % required for citecompletion, not work in mainfile
\end{document}

